% ECONOMICS_PROBLEM: {"id": "G21-547", "title": "Long-run general-equilibrium effects of microfinance", "jel_code": "G21", "source_id": "OP-0031"}
% !TeX program = xelatex
\documentclass[11pt,a4paper]{article}
\usepackage[margin=23mm]{geometry}
\usepackage{fontspec}
\setmainfont{Latin Modern Roman}
\usepackage{amsmath,amssymb,mathtools}
\usepackage{xcolor,enumitem,booktabs,longtable,array,xurl,hyperref}
\definecolor{muted}{HTML}{53636C}
\hypersetup{colorlinks=true,urlcolor=blue,linkcolor=blue}
\setlength{\parindent}{0pt}
\setlength{\parskip}{5pt}
\newcommand{\R}{\mathbb R}
\newcommand{\E}{\mathbb E}
\newcommand{\N}{\mathbb N}
\newcommand{\ind}{\mathbf 1}
\newcommand{\Var}{\operatorname{Var}}
\newcommand{\Cov}{\operatorname{Cov}}
\newcommand{\supp}{\operatorname{supp}}
\DeclareMathOperator*{\argmin}{arg\,min}
\DeclareMathOperator*{\argmax}{arg\,max}
\newcommand{\entrymeta}[3]{{\sffamily\small\color{muted}\textbf{\texttt{#1}}\quad|\quad #2\par #3\par}}
\newcommand{\Problemref}[1]{\texttt{#1}}
\newcommand{\JELref}[2]{\texttt{#2} (JEL \texttt{#1})}
\begin{document}
\section*{Long-run general-equilibrium effects of microfinance}
\textbf{Problem ID:} \texttt{G21-547}\quad \textbf{JEL:} \texttt{G21}\par
% BEGIN ECONOMICS STATEMENT
\label{problem:OP-0031}
\label{atlas:prob:D1}
\noindent\textbf{Original atlas ID:} D1 (entry 31). \textbf{Classification:} Editorial saturation, long-run treatment, and equilibrium problem.

\noindent\textbf{Original atlas question:} Does expanded credit transform productivity and occupational structure rather than merely smooth household finance?

\subsection*{Definitions and assumptions}
Fix markets $m$, nonempty finite baseline-household sets $\mathcal I_m$, households $i\in\mathcal I_m$, and horizon $H$. Treatment is an offer of a specified credit contract, including interest, liability, repayment, and access rules, rather than borrowing itself. Let $\mathbf D_m$ be the market's vector of offers. Potential outcomes $Y_{imt}(\mathbf D_m)$ may depend on other households' offers through prices, employment, and competition. Assume partial interference: arbitrary effects within a market are permitted, but no effects cross market boundaries; alternatively replace this assumption with a specified cross-market model. For saturation $s\in[0,1]$, let $\pi_s$ be an announced randomized offer-allocation rule satisfying $\mathbb E_{\pi_s}[|\mathcal I_m|^{-1}\sum_i D_{im}]=s$ in every market. Thus saturation is the expected offer share, not the borrowing share. Define outcomes separately for real consumption, business value added, enterprise entry, and occupation; assume finite expectations and consistent measurement.
\subsection*{Formal research question}
For each outcome $Y$, define the market policy response
\[
 \mu_{Y,t}(s)=\mathbb E_m\mathbb E_{\mathbf D_m\sim\pi_s}
          \left[\frac1{|\mathcal I_m|}\sum_{i\in\mathcal I_m}
                    Y_{imt}(\mathbf D_m)\right],
 \qquad \Delta_{Y,t}(s)=\mu_{Y,t}(s)-\mu_{Y,t}(0).
\]
Identify or sharply bound $\Delta_{Y,t}(s)$ for relevant saturations and $t\leq H$, especially when an entire market receives access. Determine whether output and occupational changes persist after a stated subsidy or introductory contract expires. For national expansion, construct a model of lender entry, default, capital funding, wages, and product demand, and evaluate its transition separately from these experimental market contrasts.
\subsection*{Interpretation and scope}
Individual take-up effects differ from offer effects. An encouragement instrument identifies borrower effects only with an exclusion restriction that may fail if the offer changes expectations or competing lenders' behavior. A low-saturation experiment need not identify $s=1$: positive design support or an explicit extrapolation model is required. Occupational transformation must be defined, for example as persistent changes in enterprise value added and employment shares, rather than inferred from a single positive borrowing coefficient. Attrition, enterprise closure, and household migration are outcomes to be tracked for the baseline cohort. Consumption smoothing may improve welfare without raising measured productivity, and gains for borrowers can coexist with losses for competing businesses. Welfare analysis must include repayment, lender resource costs, risk, and distributional weights.
\subsection*{Source audit and status}
[S1] is the verified Hyderabad evaluation, [S2] a six-study synthesis, and [S3] a Bayesian hierarchical analysis. These already establish context-specific and pooled evidence. They do not identify every saturation, national equilibrium, or decade-long contract design. The long-horizon scale-up estimand is editorial, not a sourced conjecture or a certified universal 2026 open problem.

\subsection*{References}
\begin{enumerate}[label=\textnormal{[S\arabic*]},leftmargin=*]
\item Abhijit Banerjee, Esther Duflo, Rachel Glennerster, and Cynthia Kinnan (2015). \emph{The Miracle of Microfinance? Evidence from a Randomized Evaluation}. American Economic Journal: Applied Economics 7(1), 22--53. \url{https://doi.org/10.1257/app.20130533}. \textit{Locator:} Primary publisher, working-paper, or author record: bibliographic information and abstract; full-text theorem verification is not claimed. \textit{Role:} Hyderabad randomized microcredit access; thematic reconstruction of the abbreviated citation.
\item Abhijit Banerjee, Dean Karlan, and Jonathan Zinman (2015). \emph{Six Randomized Evaluations of Microcredit: Introduction and Further Steps}. American Economic Journal: Applied Economics 7(1), 1--21. \url{https://doi.org/10.1257/app.20140287}. \textit{Locator:} Primary publisher, working-paper, or author record: bibliographic information and abstract; full-text theorem verification is not claimed. \textit{Role:} Synthesis of six studies; explains why program and context matter.
\item Rachael Meager (2019). \emph{Understanding the Average Impact of Microcredit Expansions: A Bayesian Hierarchical Analysis of Seven Randomized Experiments}. American Economic Journal: Applied Economics 11(1). \url{https://doi.org/10.1257/app.20170299}. \textit{Locator:} Primary publisher, working-paper, or author record: bibliographic information and abstract; full-text theorem verification is not claimed. \textit{Role:} Hierarchical analysis of treatment-effect heterogeneity; not a scale-up identification theorem.
\end{enumerate}

\subsection*{Common conventions: Source mathematical conventions}

Unless an entry states otherwise, agent and firm sets are finite. State and action spaces are measurable, strategies and outcome maps are measurable, probability laws are specified, and expectations exist for the targets under discussion. Infinite-horizon discount factors lie in $(0,1)$ and discounted objectives are finite. Norms, tolerances and complexity input sizes must be specified for computational comparisons. Constants or primitives that appear locally are inputs, not empirical facts asserted by this catalogue.

\textbf{Identification.} A statistical model $m\in\mathcal M$ implies an observable law $P_m$ and target $\theta(m)$. At law $P$, its identified set is
\[
\Theta(P;\mathcal M)=\{\theta(m):m\in\mathcal M,\ P_m=P\}.
\]
Point identification holds when this set is a singleton, or equivalently when $P_{m_1}=P_{m_2}$ implies $\theta(m_1)=\theta(m_2)$ for all compatible models. Sharp bounds mean bounds attained, or approached when the boundary is not attained, by compatible models. Writing this set is a definition, not a solution: the substantive task is to establish informative restrictions and derive or estimate the set. An unrestricted-confounding model may give only trivial bounds.

\textbf{Inference.} Identification, estimability and valid uncertainty quantification are separate questions. Fix $\alpha\in(0,1)$. A confidence procedure $C_n$ over a declared law class $\mathcal P$ has asymptotic uniform coverage at level $1-\alpha$ when
\[
\liminf_{n\to\infty}\inf_{P\in\mathcal P}\Pr_P\{\theta(P)\in C_n\}\geq1-\alpha.
\]
For set-identified targets the coverage event must be defined separately, for example coverage of the full identified set. If an entry uses identification-indexed models instead of uniquely indexed laws, the same coverage convention is applied over those models. Pointwise limits do not establish uniform validity, and asymptotic validity does not guarantee finite-sample precision. Sample sizes, dependence structures and weak-identification sequences are part of each application.

\textbf{Counterfactuals.} $Y_i(a)$ is a potential outcome under a specified intervention, including its timing, implementation and versions. It is not $Y_i$ conditional on voluntarily choosing $a$. A full intervention vector permits interference; shorthand scalar treatment notation does not silently impose no interference. Assignment, selection, attrition, anticipation and measurement must be specified. A claim about an instrument's local effect is not automatically a population policy effect.

\textbf{Economic equilibrium.} A counterfactual model includes best responses, constraints, feasibility, market clearing, state transitions and expectations consistent with the stated information. When several equilibria exist, either a declared selector is an input or the target is set-valued. Comparisons across policy regimes must use the stated selection convention. Reduced-form relationships are not presumed invariant after an equilibrium-changing intervention.

\textbf{Optimization and welfare.} Optimization is over the stated feasible set. If attainment is not established, $\sup$ or $\inf$ and $\varepsilon$-optimal choices replace an asserted maximizer or minimizer. For example, with value $V=\sup_{a\in\mathcal A}W(a)<\infty$, an $\varepsilon$-optimal set is $\{a\in\mathcal A:W(a)\geq V-\varepsilon\}$ for $\varepsilon>0$. Benefits, costs, risk and health or environmental outcomes can be aggregated only using declared units and conversion weights. Interpersonal utility weights, the treatment of fiscal transfers and foreign profits, discounting, and the behavioral welfare criterion are explicit inputs. A comparative-static derivative additionally requires existence and appropriate differentiability of the selected counterfactual.

\textbf{Sources.} Local labels [S1], [S2], and so forth restart in every question. Locators identify the relevant abstract, model, section or result at the level actually checked; a bibliographic or abstract-level verification is not represented as inspection of an entire proof. Working papers are distinguished from later publications and changing versions. Source summaries are paraphrased. All URL checks and editorial formulations refer to the audit date stated above.
% END ECONOMICS STATEMENT
\end{document}
